% part: intuitionistic-logic
% chapter: propositions-as-types

\documentclass[../../../include/open-logic-chapter]{subfiles}

\begin{document}

\olchapter{int}{pty}{Proposições como tipos}

\begin{editorial}
  Este é um rascunho \emph{muito experimental} de um capítulo sobre a
  correspondência Curry--Howard. Ele precisa de mais explicação e
  motivação, e provavelmente contém erros e omissões. A prova
  da normalização deveria ser revista e ampliada. Não há
  exemplos para o tipo produto. Conversões por permutação e simplificação
  não são abordadas. O texto fará muito mais sentido quando também
  houver material sobre o cálculo lambda tipado, que é
  basicamente pressuposto aqui. Use-o com extrema cautela.
\end{editorial}

\olimport{introduction}
%\olimport{sequent-natural-deduction}
\olimport{proof-terms}
\olimport{proofs-to-terms}
\olimport{terms-to-proofs}
\olimport{types}
\olimport{reduction}
\olimport{type-preservation}
\olimport{normalization}

\OLEndChapterHook

\end{document}
